\documentclass{article}
\usepackage{bm}
\usepackage{braket}
\usepackage{lipsum}
\usepackage{mathtools}
\usepackage{amsthm}
\usepackage{fancyhdr}
\usepackage{float}
\usepackage{listings}
\lstset{breaklines=true}
\usepackage{graphicx}
\usepackage{xcolor}
\usepackage{titlesec}
\usepackage{titling}
\usepackage{tabularx}
\usepackage{enumitem}
\usepackage{amsmath}
\usepackage{amsfonts}
\usepackage{hyperref}
\usepackage[margin=1in]{geometry}

\hypersetup{
	colorlinks=true,
	linkcolor=blue,
	filecolor=magenta,
	urlcolor=blue,
	citecolor=blue,
	anchorcolor=blue
}

\renewcommand{\thesubsubsection}{\Roman{subsubsection}}
\title{Math 3IA3 - Assignment 1}
\author{Matthew Farah, \href{mailto:farahm11@mcmaster.ca}{$\langle$farahm11@mcmaster.ca$\rangle$}, 400468251}
\date{\today}

\makeatletter

\pagestyle{fancy}
\fancyhead{}
\fancyhead[L]{Matthew Farah, $\langle$\href{mailto:farahm11@mcmaster.ca}{farahm11@mcmaster.ca}$\rangle$, 400468251}
\fancyhead[R]{Math 3IA3 - Assignment 1}

\newcolumntype{Y}{>{\centering\arraybackslash}X}
\setlength{\parindent}{0cm}

\setcounter{section}{1}

\newcounter{chapter}
\setcounter{chapter}{1}

\newcounter{exercise}
\renewcommand{\theexercise}{\arabic{exercise}}

\newenvironment{exercise}[1][]{
	\newpage
	\ifx\\#1\\
		\refstepcounter{exercise}
	\else
		\setcounter{exercise}{#1}
	\fi
	\addcontentsline{toc}{section}{\protect{\large{Exercise \thesection.\thechapter.\theexercise}}}
	\par
	\large\textbf{Exercise \thesection.\thechapter.\theexercise}
	\vspace{.2cm}
	\par
	\setcounter{subexercise}{0}
	\setcounter{step}{0}
}{\vspace{.5em}}

\newenvironment{solution}{
	\vspace{.5em}\par\large\textbf{{Solution:}}\par\vspace{.5em}
}{
	\vspace{.5em}
}

\newcounter{subexercise}
\renewcommand{\thesubexercise}{\alph{subexercise}}

\newenvironment{subexercise}{
	\refstepcounter{subexercise}
	\vspace{1em}\par\noindent\large\textbf{(\thesubexercise)}
	\setcounter{step}{0}
}{
	\par
	\vspace{1em}
}

\makeatother

\makeatletter

\newcounter{step}
\renewcommand{\thestep}{\Roman{step}}

\newenvironment{step}{
	\refstepcounter{step}
	\vspace{1.5em}\par\noindent\large\textbf{(\thestep)}
}{
	\par
	\vspace{1.5em}
}

\makeatother

\begin{document}

\maketitle
\pagestyle{fancy}

\tableofcontents
\newpage

\begin{exercise}[15]
	Show that if $f: A \rightarrow B$ and $G$ and $H$ are both subsets of $B$, then:
\end{exercise}

\begin{subexercise}
	$f^{-1}(G \cup H) = f^{-1}(G) \cup f{-1}(H)$.
\end{subexercise}

\begin{solution}
	We will prove that $f^{-1}(G \cup H) = f^{-1}(G) \cup f^{-1}(H)$ by first showing that $f^{-1}(G) \cup f^{-1}(H) \subseteq f^{-1}(G \cup H)$ and that $f^{-1}(G \cup H) \subseteq f^{-1}(G) \cup f^{-1}(H)$.

	\begin{step}
		Show that $f^{-1}(G) \cup f^{-1}(H) \subseteq f^{-1}(G \cup H)$.
	\end{step}

	Let $b \in (f^{-1}(G) \cup f^{-1}(H))$. Then, by definition of inverse, we know $f(b) \in G \vee f(b) \in H$. Thus, $f(b) \in G \cup H$. Taking the inverse of both sides, yields $b \in f^{-1}(G \cup H)$. Thus, since $b \in (f^{-1}(G) \cup f^{-1}(H)) \implies b \in f^{-1}(G \cup H)$, $(f^{-1}(G) \cup f^{-1}(H)) \subseteq f^{-1}(G \cup H)$.

	\begin{step}
		Show that $f^{-1}(G \cup H) \subseteq f^{-1}(G) \cup f^{-1}(H)$.
	\end{step}

	Let $b \in f^{-1}(G \cup H)$. Applying $f$ to both sides yields $f(b) \in (G \cup H) \implies f(b) \in G \vee f(b) \in H$. Thus, reapplying the inverse we have that $b \in f^{-1}(G) \vee \cup f^{-1}(H)$, which, by definition of inverse implies that $b \in f^{-1}(G) \cup f^{-1}(H)$. Thus, since $b \in f^{-1}(G \cup H) \implies b \in (f^{-1}(G) \cup f^{-1}(H))$, $f^{-1}(G \cup H) \subseteq (f^{-1}(G) \cup f^{-1}(H))$. \\

	Therefore, since $f^{-1}(G) \cup f^{-1}(H) \subseteq f^{-1}(G \cup H)$ and $f^{-1}(G \cup H) \subseteq f^{-1}(G) \cup f^{-1}(H)$, we necessarily have that $f^{-1}(G \cup H) = f^{-1}(G) \cup f^{-1}(H)$.
\end{solution}

\begin{subexercise}
	$f^{-1}(G \cap H) = f^{-1}(G) \cap f^{-1}(H)$.
\end{subexercise}

\begin{solution}
	We will likewise prove that $f^{-1}(G \cap H) = f^{-1}(G) \cap f^{-1}(H)$ by first showing that $f^{-1}(G \cap H) \subseteq (f^{-1}(G) \cap f^{-1}(H))$ and that $(f^{-1}(G) \cap f^{-1}(H)) \subseteq f^{-1}(G \cap H)$.

	\begin{step}
		Show that $f^{-1}(G \cap H) \subseteq (f^{-1}(G) \cap f^{-1}(H))$.
	\end{step}

	Let $b \in f^{-1}(G \cap H)$. Therefore, applying the function to both sides, we know that $f(b) \in (G \cap H) \implies f(b) \in G \wedge f(b) \in H$. Thus, applying the inverse, we have $b \in f^{-1}(G) \wedge b \in f^{-1}(H)$ and thus, by definition of intersection, we have $b \in f^{-1}(G) \cap f^{-1}(H)$. Thus, since $b \in f^{-1}(G \cap H) \implies b \in (f^{-1}(G) \cap f^{-1}(H))$, $f^{-1}(G \cap H) \subseteq (f^{-1}(G) \cap f^{-1}(H))$.

	\begin{step}
		$(f^{-1}(G) \cap f^{-1}(H)) \subseteq f^{-1}(G \cap H)$.
	\end{step}

	Let $b \in f^{-1}(G) \cap f^{-1}(H)$. Therefore, applying the function to both sides, we know that $f(b) \in G \wedge f(b) \in H$. Thus, by definition of inverse, we know that $f(b) \in (G \cap H)$, and by reapplying the inverse to both sides, we yield $b \in f^{-1}(G \cap H)$. Thus, since $b \in (f^{-1}(G) \cap f^{-1}(H)) \implies b \in f^{-1}(G \cap H)$, $(f^{-1}(G) \cap f^{-1}(H)) \subseteq f^{-1}(G \cap H)$. \\

	Therefore, since $f^{-1}(G \cap H) \subseteq (f^{-1}(G) \cap f^{-1}(H))$ and $(f^{-1}(G) \cap f^{-1}(H)) \subseteq f^{-1}(G \cap H)$ we necessarily have that $(f^{-1}(G) \cap f^{-1}(H)) = f^{-1}(G \cap H)$

\end{solution}

\setcounter{chapter}{2}

\begin{exercise}
	Conjecture a formula for the sum of the first $n$ odd natural numbers $1 + 3 + \ldots + (2n - 1)$ and prove your formula using mathematical induction.
\end{exercise}

\begin{solution}

	\vspace{.75cm}

	\begin{center}
		\begin{tabularx}{\linewidth}{X | X}
			First $n$ odd natural numbers & Sum of first $n$ odd natural numbers \\
			\hline
			1 & 1 = 1 \\
			2 & 1 + 3 = 4 \\
			3 & 1 + 3 + 5 = 9 \\
			4 & 1 + 3 + 5 + 7 = 16 \\
			5 & 1 + 3 + 5 + 7 + 9 = 25 \\
		\end{tabularx}
	\end{center}

	\vspace{.75cm}

	This table leads me to conjecture that the sum of the first $n$ odd natural numbers is equal to $n^2$. To attempt to prove this, O will use mathematical induction. Let $P(n)$ be the statement $P(n): 1 + 3 + \ldots + (2n - 1) = n^2$.

	\begin{step}
		Prove the base case (Show $P(1)$ is true).
	\end{step}

	\begin{align*}
		&\text{LHS:} & \text{RHS:} \\
		& = 1 & = 1^2 \\
		& = 1 & = 1 \\
	\end{align*}

	Therefore, since LHS $=$ RHS, $P(1)$ holds.

	\begin{step}
		State the induction hypothesis (State $P(k)$).
	\end{step}

	Let $k \in \mathbb{N}$ be arbitrary. We assume, for the purposes of our induction hypothesis that the statement:

	\begin{align*}
		P(k): 1 + 2 + \ldots + (2k - 1) &= k^2
	\end{align*}

	is true.

	\begin{step}
		Prove the induction step (Prove $P(k) \implies P(k + 1)$).
	\end{step}

	\begin{align*}
		1 + 2 + \ldots + (2k - 1) + (2k - 1) &= k^2 + (2k + 1) \\
		&= (k + 1)^2 \\
	\end{align*}

	Thus, since $P(1)$ holds true, and for arbitrary $k \in \mathbb{N}$, $P(k) \implies P(k + 1)$, by the principle of mathematical induction, $P(n)$ holds $\forall n \in \mathbb{N}$.
\end{solution}

\setcounter{section}{2}

\setcounter{chapter}{1}

\begin{exercise}[6]
	Use the argument in the proof of Theorem 2.1.4 to show that there does not exist a rational number $s$ such that $s^2 = 6$.
\end{exercise}

\begin{solution}
	Let's assume, for the purposes of deriving a contradiction, that such a rational number $s$ exists. We can therefore express $s$ in lowest form $\frac{p}{q}$ where $p$ and $q$ are coprime. We can therefore surmise that:

	\begin{align*}
		\left(\frac{p}{q} \right)^2 &= 6 \\
		p^2 &= 6q^2 \\
	\end{align*}

	Therefore, we can express $p^2$ as a multiple of $6q^2$ and thus $p^2$ must be an even number, since 6 is trivially divisible by 2. Since $p^2$ is even, we therefore know $p$ must be even as well. Thus, we can express $p$ as $2k$, where $k \in \mathbb{Z}$. Following this, we find that:

	\begin{align*}
		\frac{p^2}{q^2} &= 6 \\
		\frac{(2k)^2}{q^2} &= 6 \\
		\frac{4k^2}{q^2} &= 6 \\
		4k^2 &= 6q^2 \\
		2k^2 &= 3q^2 \\
	\end{align*}

	Since $2k^2$ is necessarily even, this implies that $3q^2$ must also be even. However, for $3q^2$ to be even, that must mean that $q^2$ is even since $3$ is not an even number. Thus, if $q^2$ is even then so must $q$. However, this presents a contradiction to our earlier assumption that $p$ and $q$ are coprime, since we have just deduced that both $p$ and $q$ are both divisible by 2. Therefore, by contradiction, there cannot exist a rational number $s$ such that $s^2 = 6$.
\end{solution}

\setcounter{chapter}{3}

\begin{exercise}[6]
	Let $S$ be a nonempty subset of $\mathbb{R}$ that is bounded below. Prove that $\inf{S} = \sup{\set{-s : s \in S}}$.
\end{exercise}

\begin{solution}
	\begin{proof}
		Appealing to the completeness property of $\mathbb{R}$, we can assert that $\inf{S}$ must exist since it is bounded below. Thus, let $m = \inf{S}$. By definition of the infimum, we know that $\forall s \in S$, $m \le s$ and therefore $-m \ge -s$. Therefore, we know that $-m$ must upperbound the set $\set{-s : s \in S}$.

		To prove that $-m$ must be the least upperbound, and thus the supremum, of $\set{-s : s \in S}$, consider that for arbitrary $\epsilon > 0$, $\exists s_0 \in S$ such that $m + \epsilon > s_0$. Therefore, $-m -\epsilon < s_0$, meaning, since $\epsilon > 0$ was arbitary that $-m = \sup{\set{-s : s \in S}}$. Therefore, deferring to out priot definition of $m$, we can see that:

		\begin{align*}
			m &= m \\
			m &= -(-m) \\
			\inf{S} &= -(\sup{\set{-s : s \in S}}) \\
			\inf{S} &= -\sup{\set{-s : s \in S}} \\
		\end{align*}
	\end{proof}
\end{solution}

\setcounter{chapter}{5}

\begin{exercise}[12]
	Give the two binary representations of:
\end{exercise}

\begin{subexercise}
	$\frac{3}{8}$
\end{subexercise}

\begin{solution}
	Since $\frac{3}{8} < \frac{1}{2}$, we take $a_1$ to be equal to 0 and therefore we have:

	\begin{gather*}
		\frac{3}{8} \in \left[ 0, \frac{1}{2} \right] \\
	\end{gather*}

	Now, since $\frac{3}{8} > \frac{1}{4}$, we can take $a_2$ to be equal to 1, thus giving:

	\begin{gather*}
		\frac{3}{8} \in \left[ \frac{1}{4}, \frac{1}{2} \right] \\
	\end{gather*}

	However, now the point of bisection of the interval $\left[ \frac{1}{4}, \frac{1}{2} \right]$ is given by $\frac{1}{4} + \frac{\frac{1}{2} - \frac{1}{4}}{2} = \frac{3}{8}$. Therefore, we can express $\frac{3}{8}$ in binary form as:

	\begin{gather*}
		\frac{3}{8} = \left( 0.011000 \ldots \right)_2 \\
		\text{or} \\
		\frac{3}{8} = \left( 0.010111 \ldots \right)_2 \\
	\end{gather*}
\end{solution}

\begin{subexercise}
	$\frac{7}{16}$
\end{subexercise}

\begin{solution}
	Since $\frac{7}{16} < \frac{1}{2}$, we take $a_1$ to be equal to 0 and therefore we have:

	\begin{gather*}
		\frac{7}{16} \in \left[ 0, \frac{1}{2} \right] \\
	\end{gather*}

	Now, since $\frac{7}{16} > \frac{1}{4}$, we can take $a_2$ to be equal to 1, thus giving:

	\begin{gather*}
		\frac{7}{16} \in \left[ \frac{1}{4}, \frac{1}{2} \right] \\
	\end{gather*}

	Since $\frac{7}{16} > \frac{3}{8}$, we set $a_3$ equal to 1 and attain:

	\begin{gather*}
		\frac{7}{16} \in \left[ \frac{3}{8}, \frac{1}{2} \right]
	\end{gather*}

	However, now the point of bisection of the interval $\left[ \frac{3}{8}, \frac{1}{2} \right]$ is given by $\frac{3}{8} + \frac{\frac{1}{2} - \frac{3}{8}}{2} = \frac{7}{16}$. Therefore, we can express $\frac{7}{16}$ in binary form as:

	\begin{gather*}
		\frac{7}{16} = \left( 0.0111000 \ldots \right)_2 \\
		\text{or} \\
		\frac{7}{16} = \left( 0.0110111 \ldots \right)_2 \\
	\end{gather*}

\end{solution}









\end{document}
